\section{The cotangent direct images and the decisive period scalar}
\label{sec:partA}

\subsection{A.1. Local logarithmic bundles and a duality calculation}

We use the marked threefold of the question, with its arbitrary admissible
constant $c$.  Coefficients of coherent cohomology are complex.  In this
section $\omega_f=K_X\otimes f^*K_B^{-1}$,
$\mathcal L=\Omega_B^1(\log\Sigma)=\OO_B(1)$,
$\mathcal A=\Omega_X^1(\log D)$, and
$\mathcal E=\mathcal A/f^*\mathcal L$.  We write $\gamma,u$ also for
their images in $R^1f_*\OO_X$; this convention concerns cohomology
classes and does not identify them with differential-form representatives.

\begin{claim}{A1.0}{Proved}{The structure-sheaf input in the precise form used below}
The direct images $R^qf_*\OO_X$, for $q=0,1,2$, are
$\OO$, $G=\OO\gamma\oplus\OO(-1)$, and $\OO(-1)$.
The inclusion $\OO\gamma\subset G$ is a subbundle.  At the cusp,
$(\gamma,u)$ is a holomorphic frame, and cup product identifies
$\bigwedge^2G$ with $R^2f_*\OO_X$.  In the cyclic covering coordinate
$s$, the action on $G$ has characters $(1,\zeta^{-a})$, where
$(m,a)=(3,2)$ or $(4,1)$ and $\zeta=e^{-2\pi i/m}$.
These assertions concern full scheme fibres, including $3S_1$ and $4S_2$.
\end{claim}
\begin{proof}
We derive the assertion from the local charts and the discriminant;
the same direct images are stated in Proposition 9.13 of
\citet{AC2026}, but that proposition is not an input to this proof.
Flatness holds
because a nonzero local parameter on the smooth base is a nonzero divisor
in the local ring of the smooth total space.  The normalization complex
of the cusp has rational surface and curve strata and the cellular
cochain complex of a torus as its constant part.  It gives
$h^q(\OO_W)=(1,2,1)$; the two degree-one classes $\gamma,u$ have nonzero
cup product.  To compute the determinant line without assuming a
canonical-bundle formula, we first retain the two finite stabilizers.
Let $\mathcal B$ be the root curve of orders three and four over $B$,
with root lines $R_1,R_2$, and put $R=R_1R_2^{-1}$.  The presentation
$3R_1=4R_2=H$ gives $\operatorname{Pic}(\mathcal B)=\Z R$,
$R_1=R^4$, $R_2=R^3$ and $\OO_{\mathcal B}(p_0)=R^{12}$;
one may also verify this by removing a line bundle's two stabilizer
characters and descending it to $\Pone$.

Write $\bar f:X\to\mathcal B$.  Translation-invariant vertical
fields on the torus charts, extended as logarithmic fields on the cusp
charts, have transition matrices depending only on the base.  They
give a rank-two bundle $V$ on $\mathcal B$ whose pullback is the
relative logarithmic tangent bundle.  Its fixed field $\xi$ is a
nowhere-zero subline, including in the logarithmic frame at the cusp.
The quotient line $Q=V/\OO\xi$ has automorphy factors $-1/\tau$
and $1/\tau$ for the two finite generators, read directly from the
linear part of the given affine deck transformations.  Thus
$\Delta(\tau)^{-1}$ is a meromorphic section of $Q^{12}$.
The discriminant product has no zero in the upper half-plane and
has a simple zero in the cusp coordinate, since
$\tau=(\log t_c)/(2\pi i)+h(t_c)$ with $h$ holomorphic.  Its divisor
therefore proves $Q^{12}=\OO_{\mathcal B}(-p_0)=R^{-12}$.
The displayed torsion-free Picard group gives
$\det V=Q=R^{-1}$.

The exact logarithmic tangent sequence has quotient
$\bar f^*T_{\mathcal B}(-p_0)$.  Its determinant, together with
$W=\bar f^*p_0$, gives
$\omega_{X/\mathcal B}=\bar f^*(\det V)^{-1}$.
The geometric fibres on the root charts have cohomology dimensions
$(1,2,1)$: they are smooth tori at the two finite points, and the
normalization calculation above gives the same result at the cusp.
Analytic base change makes the three root direct images locally free.
Connectedness gives $\bar f_*\OO_X=\OO_{\mathcal B}$, and fibrewise
Gorenstein duality with its trace pairing now yields
$R^2\bar f_*\OO_X=\det V=R^{-1}$.  The invariant local exponents
of $R^{-1}$ are two at the order-three point and one at the
order-four point.  Removing these characters leaves
$R^{-1}R_1^{-2}R_2^{-1}=R^{-12}$, so its coarse direct image is
$\OO_B(-1)$.  This supplies the degree used below independently of
the asserted direct-image or canonical-bundle formulas in \citet{AC2026}.

The same base-change calculation identifies $\gamma,u$ with a
frame through the cusp.  At a multiple fibre finite pushforward and the
Reynolds projector identify cohomology with cyclic invariants.  If $b_0$
and $b_1$ have characters $1,\zeta^{-a}$, their invariant frames are
$b_0,s^ab_1$ and $s^a(b_0\wedge b_1)$.  Their products have exactly
these exponents, so the determinant comparison is an isomorphism,
including after reduction modulo $s^m$.  Thus $\deg G=-1$.
The invariant flat class $\gamma$ does not enter the holomorphic
two-plane spanned by
$\sigma_1=6\mu\gamma+\tau u+w$ and
$\sigma_2=\beta\gamma+\mu u+\delta$: comparison of the $w,\delta$
coefficients proves this even on the smooth cyclic covers at their
centres.  It therefore gives the stated subbundle everywhere.
The quotient is $\OO(-1)$ and the extension splits because
$H^1(\OO(1))=0$.  This also specifies which subline in the otherwise
noncanonical splitting is used below.  Coherent Leray now gives
$(h^q(\OO_X))_{q=0}^3=(1,1,0,0)$, since
$H^0(\OO(-1))=H^1(\OO(-1))=H^1(\OO)=0$.
\end{proof}

\begin{claim}{A1.1}{Proved}{Local freeness and the determinant line}
Both $\mathcal A$ and $\mathcal E$ are locally free, of ranks three and
two.  The vertical vector $\xi=\partial/\partial\zeta_2$ is nowhere
zero as a section of $\mathcal E^\vee$.  There is an exact sequence
of holomorphic vector bundles
\begin{equation}\label{eq:A-E-line}
0\longrightarrow N\longrightarrow\mathcal E
 \xrightarrow{\langle\cdot,\xi\rangle}\OO_X\longrightarrow0,
\qquad N=\det\mathcal E=\omega_f(-2S_1-3S_2).
\end{equation}
\end{claim}
\begin{proof}
At a smooth point of a reduced cusp branch the assertion is the usual
coordinate calculation.  At a double point write $t_c=xy$ with third
coordinate $z$; then
$\mathcal A=R\,d\log x\oplus R\,d\log y\oplus R\,dz$ and its
base subline is $R(d\log x+d\log y)$.  At either triple point write
$t_c=xyz$; then $\mathcal A$ has the three logarithmic coordinate
basis vectors and its base subline is their sum.  In each case that
subline is a direct summand.  These calculations apply on every one
of the three double curves, including their endpoints; no simplicity
of the global divisor is assumed.

At $S_j$ we use the actual free cyclic quotient
$q:J'_j\to N_j$, which is finite and etale on total spaces, and
$t_j=s^{m_j}$ on the base.  On $J'_j$ the logarithmic conormal is
$m_j\,ds/s$, and its coefficient is a unit over $\C$.  It is again a
direct summand of the logarithmic cotangent bundle.  The comparison
$q^*\mathcal A=\Omega_{J'_j}^1(\log(s=0))$ identifies the quotient
bundles and their exact sequence; faithful flat descent proves the
claim downstairs.  This is not a replacement of $m_jS_j$ by $S_j$.

On the marked torus family $\xi$ is the constant second vertical
vector.  The linear monodromies fix it, and translations preserve it.
On a cusp toric chart it is a nonzero constant vector in the relative
logarithmic tangent lattice of rank two.  A toric coordinate change
$x_i\mapsto c_i(t_c)t_c^{b_i}x_i$ changes $d\log x_i$ by a multiple
of $d\log t_c$, so this constant relative vector descends and is
nowhere zero in that bundle.  It can vanish as an ordinary tangent
vector at a toric stratum; this does not imply vanishing in the
logarithmic bundle.  Hence evaluation is surjective and its kernel
is the determinant line.  Taking determinants of
$0\to f^*\mathcal L\to\mathcal A\to\mathcal E\to0$ gives
\[
\det\mathcal E=K_X(D)\otimes f^*\mathcal L^{-1}
=\omega_f\bigl(D-f^*\Sigma\bigr)
=\omega_f(-2S_1-3S_2),
\]
which proves \eqref{eq:A-E-line}.
\end{proof}

\begin{claim}{A1.2}{Proved}{The logarithmic direct images by vertical twisting and duality}
One has, as holomorphic vector bundles on $B$,
\[
\begin{array}{c|ccc}
q&0&1&2\\\hline
R^qf_*\mathcal E&\OO\oplus\OO(-1)&
 \OO^2\oplus\OO(-1)^2&\OO\oplus\OO(-1).
\end{array}
\]
They commute with base change to every full scheme fibre.
\end{claim}
\begin{proof}
Put $H=2S_1+3S_2$.  We first compute the line bundle $\OO_X(H)$,
rather than a tensor product of two separately descended lattices.
On the cover at $p_j$, its sections allow poles of order at most $m_j-1$
in $s$.  A fixed character has invariant exponent zero and gains no
pole, whereas the variable character $\zeta^{-a}$ has least allowed
invariant exponent $a-m$.  Thus the inclusion
$R^1f_*\OO\to R^1f_*\OO(H)$ is the identity on $\OO\gamma$ and
changes $s^ab_1$ to $s^{a-m}b_1=t_j^{-1}s^ab_1$ on its quotient.
It makes exactly one upper modification at each finite point.
The same computation applies to $R^2$.  No modification is made at
the cusp or a regular fibre.  The result is
\begin{equation}\label{eq:A-H-images}
f_*\OO(H)=\OO,\qquad
R^1f_*\OO(H)=\OO\oplus\OO(1),\qquad
R^2f_*\OO(H)=\OO(1).
\end{equation}
Here the middle extension splits because $H^1(\OO(-1))=0$.
The local invariant modules are free.  Tensoring them with
$\C\{t_j\}/(t_j)$ gives invariants of
$s^{-(m_j-1)}\C\{s\}/(s^m)$, with the indicated surviving character
basis vectors.  The Reynolds idempotent is linear over the base,
so this calculation also proves the comparison with the nonreduced
fibre.  At the cusp the line bundle is trivial.  Cohomology and base
change, applied successively from the top degree, proves the claimed
local freeness and comparison for \eqref{eq:A-H-images}.

For clarity the relative duality map we use is the cup product followed
by the analytic trace, giving a morphism of bundles
\[
R^qf_*(\omega_f\otimes\OO(-H))\longrightarrow
\bigl(R^{2-q}f_*\OO(H)\bigr)^\vee.
\]
Every fibre is a compact Gorenstein surface: it is the Cartier fibre
in a smooth threefold.  Serre duality on that scheme makes the map
an isomorphism on its residue field, also for $m_jS_j$.  The constant
dimensions just computed imply local freeness on both sides, so
Nakayama's lemma makes the bundle map an isomorphism.  Therefore
\begin{equation}\label{eq:A-N-images}
(f_*N,R^1f_*N,R^2f_*N)
=(\OO(-1),\OO\oplus\OO(-1),\OO).
\end{equation}
This is the explicit reduction of analytic relative duality needed
here; the scheme fibre is retained throughout.

The connecting maps of \eqref{eq:A-E-line} are zero on a regular
fibre: its relative cotangent is a trivial bundle and evaluation on
a constant nonzero vector splits.  Source and target direct images
in \eqref{eq:A-N-images} and Statement A1.0 are vector bundles,
so these holomorphic connecting maps, zero on a dense open set, are
zero everywhere.  We obtain extensions in each degree of the listed
$R^qN$ by $R^q\OO$.  Every resulting pair of line-bundle degrees is
zero or minus one; their extension groups are
$H^1(\OO(-1)),H^1(\OO),H^1(\OO(1))$, all zero.  These extensions
give exactly the three asserted bundles.  Applying the already proved
fibre comparisons to the short exact sequences proves base change
for $R^q\mathcal E$ as well.  This argument uses the vertical twist
$\OO(H)$ and duality; it does not use the workbench's filtration of
$f_*\mathcal E\otimes G$.
\end{proof}

\begin{claim}{A1.3}{Proved}{All absolute cotangent direct images}
For every admissible $c$, the six direct images are
\begin{equation}\label{eq:A-absolute-table}
\begin{array}{c|cc}
q&R^qf_*\Omega_X^1&R^qf_*\Omega_X^2\\\hline
0&\OO(-2)&\OO(-2)\oplus\OO(-1)\\
1&\OO\oplus\OO(-1)^2&\OO(-2)\oplus\OO(-1)^2\\
2&\OO\oplus\OO(-1)&\OO.
\end{array}
\end{equation}
In particular every listed sheaf is torsion-free.
\end{claim}
\begin{proof}
We use Statements A2.1, A2.2 and A2.4 below, whose proofs do not use
\eqref{eq:A-absolute-table}.  They give
$\coker\kappa_0=\C_{p_1}\oplus\C_{p_2}$,
surjectivity of $\kappa_1$, and rank one on its global sections.
Write $\mathcal K=\ker\kappa_1$.  The logarithmic conormal sequence
then gives
\[
f_*\mathcal A=\OO(1),\quad
0\to T\to R^1f_*\mathcal A\to\mathcal K\to0,\quad
R^2f_*\mathcal A=\OO\oplus\OO(-1),
\quad T=\C_{p_1}\oplus\C_{p_2}.
\]
Let $\nu:\widetilde D=\mathrm{dP}_6\sqcup S_1\sqcup S_2\to D$
be normalization.  The local bases in Statement A1.1 give the exact
residue sequence
\[
0\to\Omega_X^1\to\mathcal A\to\nu_*\OO_{\widetilde D}\to0.
\]
Its three residue direct images are
$\C_{p_0}\oplus\C_{p_1}\oplus\C_{p_2}$,
$\C_{p_1}\oplus\C_{p_2}$, and zero: the normalization of $W$ is
rational, and a bielliptic surface has $(h^0,h^1,h^2)(\OO)=(1,1,0)$.
In degree zero the map from $\OO(1)$ is local evaluation with the
nonzero multipliers $1,3,4$.  Its kernel is $\OO(-2)$.

We check the torsion cancellation, rather than discarding it.  At
$p_j$, in the invariant frame $(b_0,s^ab_1)$ of $G$, residue of
$(dt_j/t_j)(x b_0+y s^ab_1)$ in $H^1(S_j,\OO)$ is $m_jx[b_0]$.
The residue factors through the base logarithmic line on the covering
chart.  Naturality of the connecting map consequently gives
$\operatorname{res}_1\kappa_0=0$.  The determinant of $\kappa_0$
has a simple zero, so its reduction has rank one, and its image is
exactly the one-dimensional kernel $\C s^ab_1$ of this residue.
Hence $\coker\kappa_0\to H^1(S_j,\OO)$ is an isomorphism.
Both torsion summands cancel, giving the canonical identification
$R^1f_*\Omega_X^1=\mathcal K$ and the asserted $R^2$.
The rank-one global map
$\OO^2\oplus\OO(-1)^2\to\OO$ has a nonzero constant component
on a trivial summand.  Eliminating all other components by bundle
automorphisms gives
$\mathcal K=\OO\oplus\OO(-1)^2$.

Finally $\Omega_X^2=(\Omega_X^1)^\vee\otimes K_X$.  Relative
duality, with the locally free direct images now established, gives
\[
R^qf_*\Omega_X^2
=\bigl(R^{2-q}f_*\Omega_X^1\bigr)^\vee\otimes K_B.
\]
Substitution of $K_B=\OO(-2)$ gives every entry of the second column.
No Euler-characteristic identity is used to select a Hodge number.
\end{proof}

\subsection{A.2. An exact modular determinant and two numerical routes}

The scalar depends on the fourth root chosen in $t-1=s^4$.  We make
that ambiguity explicit.  We choose the local branch for which, on
the positive $t>1$ arc leading to the distinguished cusp,
$\tau=iy$ with $y>1$, and
\begin{equation}\label{eq:A-s-branch}
\tau=i+\frac{i}{\pi\sqrt{E_4(i)}}s^2+O(s^4),\qquad
\operatorname{Im}\mu_s(0)>0.
\end{equation}
The nonvanishing proved below makes the second condition meaningful;
it fixes the remaining sign of $s$.  Equivalently it is the lower
cusp-regular zero branch of the Weierstrass function in Statement A2.3.
Replacing $s$ by $i^js$ multiplies the reported derivative by $i^{-j}$.
The workbench instead normalizes $ds/dz(z_2)=1$; its number is obtained
by the chain rule and is not numerically the same unspecified coordinate.

\begin{claim}{A2.1}{Proved}{The complete raw period determinant}
On the marked upper half-plane,
\begin{equation}\label{eq:A-modular-K}
6\mu_z^2-\tau_z\beta_z
=\tau_z^2\frac{E_4(\tau)^2}{E_6(\tau)}.
\end{equation}
The right side extends holomorphically through the apparent poles.
Its exact orders are two at $z_1$ and zero at $z_2$; it has no other
zero away from their translates.  The logarithmic connecting map
$\kappa_0:f_*\mathcal E\to\mathcal L\otimes G$ has divisor
$\operatorname{divisor}(\det\kappa_0)=p_1+p_2$.
\end{claim}
\begin{proof}
Consider the holomorphic quadratic differential
$\mathscr Q=6(d\mu)^2-d\tau\,d\beta$.  Direct differentiation of
the two affine transformation laws in AC26 Definition 3.1 gives
\[
g_j^*\mathscr Q=\tau^{-2}\mathscr Q\quad(j=1,2).
\]
In particular $g_2^2$ fixes $\tau$ and sends
\[
\mu\longmapsto1-\tau-\mu,\qquad
\beta\longmapsto\beta+6\tau+12\mu-6,
\]
which leaves $\mathscr Q$ exactly unchanged.  The symbolic residuals
of all three identities are zero in \texttt{a\_calculations.py}.
The kernel of
$\Delta=\Z/3*\Z/4\to\Z/3*\Z/2=\operatorname{PSL}_2(\Z)$
is the normal subgroup generated by $g_2^2$.  The squared multiplier
of $\mathscr Q$ is the ordinary modular multiplier
$(c\tau+d)^{-2}$, so it is one on every kernel element.  Two generic
points of the marked cover with the same $\tau$ have the same $t$
and differ by an element of $\Delta$; its modular image fixes a
generic point, hence is the identity.  It follows that
$k=\mathscr Q/(d\tau)^2$ is a single-valued meromorphic function of
$\tau$, transforming as a modular form of weight two.  This kernel
argument is needed: mere equivariance of the separate period functions
would not prove descent.

At $z_1$, write a linearizing parameter $r$ with
$g_1r=\zeta r$, $\zeta=e^{-2\pi i/3}$.  If $\mathscr Q=K_rdr^2$
has first nonzero coefficient of order $n$, its transformation law
at the fixed point reads $\zeta^{n+2}=\rho^{-2}=\zeta$.
Thus $n\equiv2\pmod3$, and $K_r$ vanishes to order at least two.
The map $\tau$ is unramified there by AC26 Theorem 3.4(i).  At
$z_2$, $\tau-i$ has order two and $K_s$ is holomorphic; therefore
$k$ has at most a simple pole at $i$.  Away from these points
$\tau$ is unramified.  At the cusp $\mu$ and $\beta+\tau$ are
holomorphic in $q$, so $k=1+O(q)$.

The modular function $k E_6/E_4^2$ is therefore holomorphic at
$\rho$, at $i$, and at the cusp, as well as everywhere else.
The usual zero statements for $E_4,E_6$ (simple at $\rho,i$,
respectively) are the modular facts used in AC26 Section 3.1.
This function descends to the compact modular curve $X(1)=\Pone$,
so it is constant.  Its cusp value is one.  This proves
\eqref{eq:A-modular-K} and the exact two local orders, without the
workbench's degree count for the coarse determinant line.

For completeness we translate to the actual coarse lattices.  Choose
cover eigenframes $e_0,e_1$ of $\mathcal E$ with characters
$1,\zeta^a$, and $b_0,b_1$ of $G$ with characters
$1,\zeta^{-a}$.  Their invariant bases are
$e_0,s^{m-a}e_1$ and $b_0,s^ab_1$.  The logarithmic base frame is
$\ell=dt_j/t_j=m\,ds/s$.  Changing determinants from these frames
to the marked ones gives
\[
\det\kappa_0=\epsilon(t_j)m^{-2}s^{m+2-2a}K_s,
\qquad \epsilon(0)\ne0.
\]
This equality follows by taking the two factors $s/m$ from replacing
$ds$ by $\ell$, the factor $s^{m-a}$ in the source, and the factor
$s^{-a}$ in the target.  Thus at $(m,a)=(3,2)$ the total order in
$s$ is $1+2=3$, and at $(4,1)$ it is $4+0=4$.
Both coarse zeros are simple.  At the cusp, the frame change has
no elementary divisor and $dZ/d\tau$ tends to
$\left(\begin{smallmatrix}0&1\\-1&0\end{smallmatrix}\right)$,
with determinant one.  Formula \eqref{eq:A-modular-K} excludes any
further zero, proving the asserted divisor.
\end{proof}

\begin{claim}{A2.2}{Proved}{The evaluated scalar and the two local leading terms}
In the coordinate \eqref{eq:A-s-branch},
\begin{equation}\label{eq:A-mu-value}
\mu_s(0)=i\,C_4,
\qquad
C_4=\frac{\sqrt{2/3}}{\pi E_4(i)^{1/4}}
=\frac{4\sqrt\pi}{3^{3/4}\Gamma(1/4)^2}>0.
\end{equation}
Numerically,
\[
C_4=0.2366093140807852246877051134504496062943051340858243372192243072\ldots.
\]
For $r^3=t$, with $r>0$ on the arc $0<t<1$ and $\tau$ travelling
from $\rho$ to $i$ along the unit circle, set
\[
\lambda=\frac{3e^{5\pi i/6}}{2\pi E_6(\rho)^{1/3}}.
\]
Then $\tau=\rho+\lambda r+O(r^2)$ and
\[
K_r=-\frac{3i\lambda}{2\pi}r^2+O(r^3).
\]
The leading coefficient has real and imaginary parts
\[
\begin{aligned}
\operatorname{Re}K_2&=0.08010221108753827140763650309934451243\ldots,\\
\operatorname{Im}K_2&=0.13874109940222333968402480721280640441\ldots.
\end{aligned}
\]
\end{claim}
\begin{proof}
Ramanujan's differential relations give
$E_6'(i)=-\pi i E_4(i)^2$ and
$E_4'(\rho)=-(2\pi i/3)E_6(\rho)$.
Using $j=1728E_4^3/(E_4^3-E_6^2)$, we obtain
\[
j(\tau)-1728=-1728\pi^2E_4(i)(\tau-i)^2+O((\tau-i)^3).
\]
Since $j=1728(1+s^4)$, the chosen branch is precisely
\eqref{eq:A-s-branch}.  Formula \eqref{eq:A-modular-K} now gives
\[
K_s(0)=-\frac4{\pi^2\sqrt{E_4(i)}}=6\mu_s(0)^2,
\]
because $\tau_s(0)=0$.  This proves nonvanishing and the first
formula in \eqref{eq:A-mu-value}; the sign is fixed by the stated
coordinate convention.  Similarly
$j(\tau)=-512i\pi^3E_6(\rho)(\tau-\rho)^3+cdots$.
This gives the stated $\lambda$, and substitution into
\eqref{eq:A-modular-K} gives the coefficient $-3i\lambda/(2\pi)$.

There is also a geometric nonvanishing proof independent of this
modular determinant.  Using the section identification proved in
A2.3, set $x=s^2X,y=s^3Y,t=1+s^4$.  The smooth family becomes
\[
Y^2=X^3-3(1+s^4)^3X+2(1+s^4)^4s^2,
\qquad P=(0,-\sqrt2(1+s^4)^2s)
\]
on the chosen branch.  At $s=0$ its equation is
$Y^2=X^3-3X$ and the point is $(0,0)$.  The first variation of the
curve itself is zero.  In the local coordinate $Y$ at that point,
$dX/(2Y)=dY/(3X^2-3)$, so the Abel integral of $P$ has derivative
$\sqrt2/3$.  The marked period is $-i\Omega_R$, rather than the
positive real period $\Omega_R$: the positive imaginary branch
$\tau=iy$, $y>1$, has elliptic scale $F/s=i/\Omega_R$ at zero.
Its period has zero first derivative.  Division gives
$\mu_s(0)=i\sqrt2/(3\Omega_R)\ne0$.
Using $g_2(\Z+i\Z)=(4\pi^4/3)E_4(i)$ and
$g_2(\Omega_R(\Z+i\Z))=12$ gives
$\Omega_R=\pi E_4(i)^{1/4}/\sqrt3$, which identifies this value
with \eqref{eq:A-mu-value}.  This also explains why choosing the
positive real period with a positive real $s$ instead would select
the reciprocal branch $\operatorname{Im}\tau<1$.

Here are the numerical methods actually run.  The first evaluates
$E_4(i)=1+240\sum_{n=1}^{100}\sigma_3(n)e^{-2\pi n}$ at 100
decimal working digits and inserts it in \eqref{eq:A-mu-value}.
The second integrates
\[
\Omega_R=\frac{\sqrt2}{3^{1/4}}
\int_0^{\pi/2}(1-\tfrac12\sin^2\theta)^{-1/2}\,d\theta,
\qquad C_4=\frac{\sqrt2}{3\Omega_R}.
\]
This is the real period of $Y^2=X^3-3X$ for $dX/(2Y)$; substituting
the three roots $\sqrt3,0,-\sqrt3$ in the elementary elliptic-integral
change of variable proves the formula.  The beta-integral evaluation
of the displayed complete elliptic integral gives the gamma
expression in \eqref{eq:A-mu-value}.  The Fourier and quadrature
answers agree with absolute discrepancy $1.79\cdot10^{-102}$.
As a third check, Newton iteration on the Weierstrass Fourier series
with $j(\tau)=1728(1+s^4)$, using the two roots adjacent to
$(1-\tau)/2$, gives their symmetric derivative.  At $s=10^{-15}$ it
is
\[
i\,0.23660931408078522468770511345048081721971218\ldots,
\]
with discrepancy $3.13\cdot10^{-32}$ and maximum residual below
$2.9\cdot10^{-101}$.  The order-three coefficient was independently
evaluated using the divisor-sum expansion of $E_6$ and its Jacobi
theta-constant formula, with agreement to the working precision.
The code, all output, and precision are supplied.  These are numerical
corroborations, not interval certificates; nonvanishing and the exact
orders were proved analytically above.  For example the elementary
bound $\sigma_3(n)\le n^4$ and $e^{-2\pi}<1/500$ bounds the omitted
Fourier tail after $N$ terms by
$250(N+1)^4 500^{-(N+1)}$; roundoff is separately controlled by
the working precision and the independent calculation.
\end{proof}

\begin{claim}{A2.3}{Proved}{Identification of the elliptic logarithm and the verified zero formula}
For the Weierstrass equation
\[
y^2=x^3-3t^3(t-1)x+2t^4(t-1)^2,
\]
the section $[\mu]$ is a generator of its rational section group and
equals one of $P_*=(0,\sqrt2t^2(t-1))$ and $-P_*$.
In particular $\wp(\mu,\tau)=0$ on every regular fibre.
Along the branch in \eqref{eq:A-s-branch}, with
$x_0=1-1/t$, the exact branch is
\begin{equation}\label{eq:A-DI-branch}
\mu=\frac{1-\tau}{2}
 +\frac{i\sqrt6}{3\pi}
 \frac{x_0^{1/4}\,{}_3F_2(1/3,2/3,1;3/4,5/4;x_0)}
 {{}_2F_1(1/12,5/12;1;1-x_0)}.
\end{equation}
Consequently $\mu-(2-\tau)/3$ has exact order two at the
order-three point in $r^3=t$.
\end{claim}
\begin{proof}
We include the comparison, since a matching $j$-map alone would not
identify a section.  Write $F$ for the elliptic scale, not the
threefold projection.  It is specified by
\[
F^2=-\frac{3t(t-1)E_4(\tau)}{\pi^2E_6(\tau)},\qquad
F/t\longrightarrow i\sqrt3/\pi
\]
at the distinguished cusp.  The local orders are two at $r=0$ and
one at $s=0$; hence the square root exists on the simply connected
marked cover.  Its transformation laws are
$F(g_1z)=-F(z)/\tau(z)$ and $F(g_2z)=F(z)/\tau(z)$.
These follow first after squaring by modularity, and the signs are
fixed by those local orders.  Direct substitution in the cubic
identifies its regular fibres with $\C/(\Z+\tau\Z)$ by
\[
x=F^2\wp(\nu,\tau),\qquad y=\tfrac12F^3\wp'(\nu,\tau).
\]
The transformation laws of $\mu$ make this a single-valued section
$P$.  Its coordinates are meromorphic at the two finite special
points and at the cusp, so are rational functions on $\Pone$.

Every rational section $Q$ admits a global holomorphic elliptic
logarithm $\nu$ on the marked cover.  Indeed the exponential map
from the trivial complex line bundle to this smooth elliptic family
is a covering map with locally uniformly discrete lattice; pullback
over a simply connected base lifts the section.  The apparent poles
in its projective Weierstrass coordinates at a finite point are
meromorphic poles of finite order: the coordinates of $Q$ are
rational in $t$, while $F$ is holomorphic with the finite orders
just computed.  Multiplying the three homogeneous coordinates by
a common power of the local uniformizer and cancelling their
common zero gives holomorphic coordinates not all zero at the
centre.  Their limiting point lies on the closed projective
Weierstrass family.  This is the valuative criterion for properness
over the local analytic discrete valuation ring; it is not an
assertion that an arbitrary punctured-disc holomorphic map to a
compact space has a removable singularity.
We can choose this logarithm bounded at the distinguished cusp.
To verify this last assertion, put $v=1/t$, $X=v^2x$, $Y=v^3y$.
The equation becomes
$Y^2=X^3-3(1-v)X+2(1-v)^2$.  A section cannot meet its nodal point
$(X,Y)=(1,0)$ at $v=0$: substituting $X=1+O(v),Y=O(v)$ leaves
nonzero coefficient one in the defining equation at order $v$.
On the smooth nodal locus the coordinate
$u=\exp(2\pi i\nu)$ is a nonzero finite coordinate, as follows
directly from
\[
X=1+\frac{12u}{(u-1)^2},\qquad
Y=\frac{12\sqrt3u(u+1)}{(u-1)^3}.
\]
The normally convergent Weierstrass Fourier series extends this
coordinate to a nearby smoothing.  The inverse function theorem
therefore writes the section as a holomorphic unit $u(v)$, also
near $u=1$, the point at infinity.  A logarithm of that unit fixes
the bounded lift.  Any two such lifts differ by an integer, since
an additional nonzero multiple of $\tau$ is unbounded.

For this lift, write the integral transformation errors as
\[
\nu(g_1z)=-\nu/\tau+\alpha_1+\beta_1(\tau-1)/\tau,
\quad
\nu(g_2z)=\nu/\tau+\alpha_2-\beta_2/\tau.
\]
Bounded cusp regularity and $g_1g_2=g_0^{-1}$ imply
$\alpha_2=-\beta_1$, $\beta_2=\alpha_1+\beta_1$.
Putting $n=-\alpha_1-2\beta_1$, $a=\alpha_1+\beta_1$, we find that
$\nu-(n\mu+a)$ satisfies the homogeneous law of AC26 Lemma 3.10
and is cusp-regular.  We recall a direct verification of the degree
used here.  The function
$F_h=E_4(\tau)^2E_6(\tau)^{1/2}/\Delta(\tau)$ on the marked cover
has orders two and one at $z_1,z_2$, and has a simple pole at the
cusp.  The square root exists because its zeros on this cover have
even order.  Modularity and these local orders fix the two signs:
$F_h(g_1z)=-F_h/\tau$, $F_h(g_2z)=F_h/\tau$.
For any holomorphic solution of the same homogeneous law, the
quotient by $F_h$ is an invariant meromorphic function of $t$.
A possible pole at the order-three point would have order at most
$2/3$ in $t$, and one at the order-four point at most $1/4$; its
integral order is therefore nonnegative.  At the cusp regularity
requires this quotient to vanish.  Thus the solution sheaf is
exactly $\OO_B(-p_0)$, with no nonzero global sections.
We obtain $\nu=n\mu+a$ and hence $Q=nP$.  The two distinct
finite values of $\mu$ exclude $nP=0$ for $n\ne0$ by cusp boundedness;
there is no torsion.  This proves that $P$ generates the section group.

We identify $P_*$ with a primitive generator by an intersection
calculation.  The rational elliptic surface and its fibre types
$\mathrm{IV}^*,\mathrm{III},I_1$ are those in Engel Section 2.
At $t=1$, the blowup chart $x=(t-1)X,y=(t-1)Y$ gives
$Y^2=(t-1)X^3-3t^3X+2t^4$; $P_*$ meets its exceptional
multiplicity-one component at $(X,Y)=(0,\sqrt2)$, away from the
identity component.  At $t=0$, two successive charts
$x=t^2X,y=t^2Y$ give
$Y^2=t^2X^3-3t(t-1)X+2(t-1)^2$; $P_*$ meets the smooth exceptional
component $Y=-\sqrt2$ at $X=0$.  These are the two nonidentity
simple components of the $E_6$ configuration.  The section is
disjoint from $O$ in every chart, including infinity.

Here the height calculation can be made directly, without assuming
primitivity.  Orthogonally project $Q-O$ in
$\operatorname{NS}(S)\otimes\Q$ off $O$, the fibre $F_S$, and the
nonidentity components.  The inverse intersection matrices are
$-C_{E_6}^{-1}$ and $-1/2$.  Their diagonal entries for the two
nonidentity simple components are $-4/3$ and $-1/2$.
Here we may number the nonidentity components so that
\[
C_{E_6}=\begin{pmatrix}
2&-1&0&0&0&0\\
-1&2&-1&0&-1&0\\
0&-1&2&-1&0&0\\
0&0&-1&2&0&0\\
0&-1&0&0&2&-1\\
0&0&0&0&-1&2
\end{pmatrix},\qquad
(C_{E_6}^{-1})_{44}=(C_{E_6}^{-1})_{66}=4/3.
\]
Since $K_S=-F_S$, adjunction gives $Q^2=O^2=-1$.  Expanding the
square of this projection gives
\[
h(Q)=2+2(Q\cdot O)-c_0(Q)-c_1(Q),\quad
c_0(Q)\in\{0,4/3\},\ c_1(Q)\in\{0,1/2\}.
\]
The projection is additive: the chord-and-tangent law makes
$Q+R-(Q+R)_{\rm section}-O$ linearly equivalent to a vertical
divisor, which the projection kills.  Its negative square therefore
satisfies $h(nQ)=n^2h(Q)$.  A section meets a single simple
component, since its base parameter has order one; consequently
the listed local corrections exhaust the possibilities.
For $Q\ne O$ we have $Q\cdot O\ge0$, so $h(Q)\ge1/6$.
For $P_*$ the explicit incidences give $h(P_*)=1/6$.
Writing $P_*=nP$ now gives $1/6=n^2h(P)\ge n^2/6$, and hence
$n=\pm1$.  Inverting the displayed $E_6$ Cartan matrix and the
chord-and-tangent formulas is included in the exact computation log.
This proves $\wp(\mu,\tau)=0$.

We checked the original theorem on page 399 of
\citet{EichlerZagier1982}, as well as equation (6) of
\citet{DukeImamoglu2008}.  With the normalized $\Delta$, their
coefficient is indeed $144\pi i\sqrt6$, their constant is
$\log(5+2\sqrt6)/(2\pi i)$, and the integrand is
$(\sigma-\tau)\Delta(\sigma)/E_6(\sigma)^{3/2}$ on the vertical
path with its continued branch.  Thus the formula in Lead 3 is
correct with its indicated sign and period ambiguity.  Theorem 1
of Duke--Imamoglu gives \eqref{eq:A-DI-branch}: we take one minus
their $z_0$, rather than $z_0-\tau$, because the latter is unbounded
at the distinguished cusp.  The chosen branch approaches
$(1-i)/2$ at $\tau=i$ and has lower imaginary cusp value.
These conditions identify the original bounded lift.  In particular
its linear coefficient is
$i\sqrt6/(3\pi\,{}_2F_1(1/12,5/12;1;1))$, agreeing with
\eqref{eq:A-mu-value}.

Finally exact elliptic addition gives
\[
x(2P_*)=9t^2/8,\quad
y(2P_*)=\frac{t^2(32-5t)}{16\sqrt2},\quad
x(3P_*)=\frac{16}{81}(t-1)(5t-32).
\]
The last expression has nonzero limit $512/81$ at $t=0$.
If $\nu=\mu-(2-\tau)/3$, then $3\nu$ is the local logarithm of
$3P$ modulo its exact period.  Since $F$ has order two and
$\wp(3\nu)=(3\nu)^{-2}+O(\nu^2)$, the equality
$x(3P)=F^2\wp(3\nu)$ forces $\nu$ to have exact order two.
It also gives its leading coefficient
$\nu_2=\pm3F_2/(16\sqrt2)$, where $F_2=\lim_{r\to0}F/r^2$;
the sign is the explicitly stated elliptic-logarithm branch.
\end{proof}

\begin{claim}{A2.4}{Proved}{The global rank and Lead 2}
The map
$\kappa_1:R^1f_*\mathcal E\to\mathcal L\otimes R^2f_*\OO=\OO$
is surjective, and $H^0(\kappa_1):\C^2\to\C$ has rank one.
Let $h=a\,d\zeta_1+d\zeta_2$ be the unique global section of
$f_*\mathcal E$ with $h(\xi)=1$.  In a genuine invariant target
frame normalized as below,
\[
\kappa_1(h\gamma)(p_2)=-\tfrac14\mu_s(0)=-\tfrac i4 C_4.
\]
The rank-one condition is equivalent to
$a\tau_z+\mu_z\not\equiv0$.  Under the same direct-image and
sheaf-surjectivity hypotheses a rank-zero global map would instead
give $h^{1,2}=2$, $h^{1,1}=3$; that alternative does not occur here.
\end{claim}
\begin{proof}
Let $d_{ij}(s)$ be the cover Kodaira--Spencer coefficients in the
eigenframes of Statement A2.1 after removing $ds$.
The invariant coarse matrix and its cup-product row have the forms
\[
[\kappa_0]=\begin{pmatrix}A&t_jB\\C&D_0\end{pmatrix},\qquad
[\kappa_1]=(-C,A,-D_0,B),
\]
in the ordered basis
$e_0b_0,s^ae_0b_1,s^{m-a}e_1b_0,e_1b_1$.  Specifically
$A=sd_{00}/m$, $B=s^{1-a}d_{01}/m$,
$C=s^{1-a}d_{10}/m$, $D_0=s^{m+1-2a}d_{11}/m$.
These are holomorphic invariant functions; the row follows by
wedging the connecting class with $b_0,b_1$.
If the row had no unit at a finite point, every entry would vanish
modulo $t_j$ and $\det\kappa_0$ would vanish to order at least two,
contradicting A2.1.  Away from the finite points, including the cusp,
$\kappa_0$ is invertible, which gives a unit in the cup-product row.
Thus $\kappa_1$ is surjective.

The lift $h$ exists and is unique because $f_*\mathcal E$ is an
extension of $\OO$ by $\OO(-1)$.  Its invariance under $g_2$ gives
$a(g_2z)=\tau a+\mu$, hence $a(z_2)=1/2$.
In the quotient frame $(\gamma,u)$ of $G$ the connecting map is
\[
\kappa_0(h)=\frac{s}{4}
\bigl((6a\mu_s+\beta_s)\gamma+(a\tau_s+\mu_s)u\bigr)\ell.
\]
Choose the variable eigenframe $b_1$ by Reynolds averaging so that
$\gamma\wedge b_1=v(s)\gamma\wedge u$, $v(0)=1$.
Then $s\,v(s)(\gamma\wedge u)\ell$ is the actual invariant
target frame; the uncorrected expression with $v=1$ need not be
invariant away from zero.  Cup product with $\gamma$ gives at the
central point $-(a\tau_s+\mu_s)/4=-\mu_s/4$, proving rank one.

To check the stated equivalence, there is a second global section
\[
q_E=\sigma_1\otimes u+6\sigma_2\otimes\gamma
\quad\text{of }R^1f_*\mathcal E.
\]
It is the negative of the relative $(1,1)$ image of the invariant
flat class $\eta=u\wedge w+6\gamma\wedge\delta$.
It extends at the cusp because the relative logarithmic forms and
$\gamma,u$ extend there.  At a finite point it is an invariant
holomorphic tensor on the smooth cover, hence belongs to the four
explicit invariant generators above.  Its cup-product row is zero,
since the two terms are $6\mu_z$ and $-6\mu_z$.
It is independent of $h\gamma$ on every regular fibre (its
$\sigma_1\otimes u$ coefficient is nonzero), so these two sections
form a basis of the two-dimensional global section space.
Consequently the global map has rank one precisely when its value
on $h\gamma$, equivalently $a\tau_z+\mu_z$, is not identically
zero.  If its rank were zero, the long exact cohomology sequence of
$0\to\mathcal K\to\OO^2\oplus\OO(-1)^2\to\OO\to0$
would give $h^0(\mathcal K)=2$, $h^1(\mathcal K)=1$.
The residue calculation in A1.3 and Leray would then give exactly
the two counterfactual Hodge numbers stated in the claim.
\end{proof}

\subsection{A.3. The Hodge table and its Leray generators}

\begin{claim}{A3.1}{Proved}{The complete Hodge table for every admissible parameter}
For every $\operatorname{Im}c<-M$, the Hodge numbers are
\[
\begin{array}{c|rrrr}
q\backslash p&0&1&2&3\\\hline
3&0&0&0&1\\
2&0&1&2&1\\
1&1&2&1&0\\
0&1&0&0&0.
\end{array}
\]
\end{claim}
\begin{proof}
Since $B$ is a curve, coherent Leray has no differential and gives
for $F=\Omega_X^p$ the exact sequences
$0\to H^1(B,R^{q-1}f_*F)\to H^q(X,F)\to H^0(B,R^qf_*F)\to0$.
Substitution of \eqref{eq:A-absolute-table}, with
$h^0(\OO)=h^1(\OO(-2))=1$ and all cohomology of $\OO(-1)$ zero,
gives columns $p=1,2$ directly.  Statement A1.0 gives column zero,
and Serre duality gives column three.  Thus the numbers are outputs
of the sheaf computation.  Replacing $\beta$ by $\beta+c$ leaves
every derivative, finite eigencharacter, invariant lattice and
logarithmic residue used above unchanged.  The toric gluing units
change, but their relative logarithmic derivatives still differ only
by base logarithmic forms.  Hence the same computations apply at
each admissible $c$, including every point arbitrarily close to the
strict boundary; no manifold at the excluded boundary is asserted.
\end{proof}

\begin{claim}{A3.2}{Proved}{Specified associated-graded generators and a global Dolbeault lift}
The Leray associated-graded spaces for the four middle groups are
\[
\begin{array}{c|cc}
&E_2^{1,q-1}&E_2^{0,q}\\\hline
H^{1,1}&H^1(\OO(-2))=\C\alpha&\C\eta_{\rm rel}\\
H^{1,2}&0&\C\upsilon\\
H^{2,1}&\C\upsilon^\vee&0\\
H^{2,2}&\C\eta_{\rm rel}^\vee&\C\alpha^\vee.
\end{array}
\]
Here $\alpha$ is the Cech class of $dt/t$ on the two standard affine
base charts, $\eta_{\rm rel}=-q_E$ is the invariant relative class
from A2.4, and $\upsilon$ is the image of the trace-dual of
$1\in f_*\OO(H)$ under $R^2f_*N\to R^2f_*\mathcal E$ and the
isomorphism $R^2f_*\Omega_X^1=R^2f_*\mathcal E$.
The superscript $\vee$ denotes the bases dual under the relative
trace and the base Serre pairing.  There exists a Dolbeault class
$Q\in H^1(X,\Omega_X^1)$ whose restriction to every regular fibre
is $\eta=u\wedge w+6\gamma\wedge\delta$.
\end{claim}
\begin{proof}
The isomorphism $f_*\Omega_X^1=\Omega_B^1$ in the residue sequence
is the actual pullback map, so $\alpha$ gives its stated Leray
class; it can be represented by a nonzero normalization of
$f^*\omega_{\rm FS}$ in Dolbeault cohomology.  The section
$\eta_{\rm rel}$ belongs to the canonical kernel identification
$R^1f_*\Omega_X^1=\ker\kappa_1$ and spans its global sections.
Leray is surjective onto this space, so any preimage defines $Q$
with the asserted fibre restriction.  Its ambiguity is exactly
$\C\alpha$; the statement does not claim a canonical global form.
The injection $R^2f_*N=\OO\to R^2f_*\mathcal E$ from
\eqref{eq:A-E-line} spans its global sections, since the quotient
is $\OO(-1)$.  This defines $\upsilon$ with no choice of a negative
degree summand.  Duality used in A1.3 gives the remaining entries
and their specified dual pairings.  Choosing a lift of
$\eta_{\rm rel}$ and of $\alpha^\vee$ gives bases of the whole
middle cohomology groups.  The associated-graded descriptions are
canonical with the named choices; the splitting of Leray vector
spaces is not asserted canonical.
\end{proof}

\subsection{A.4. Ordinary relative differentials and the correction at the cusp}

\begin{claim}{A4.1}{Proved}{All local modules needed to distinguish degrees one and two}
Let $R=\C\{x,y,z\}$.  For $t=xy$,
\[
\Omega_{R/\C\{t\}}^1=R^3/R(y,x,0)
\simeq(x,y)\oplus R,
\]
and
\[
\Omega_{R/\C\{t\}}^2\simeq
(x,y)\oplus R/(x,y),\qquad
\Omega_{R/\C\{t\}}^3\simeq R/(x,y).
\]
For $t=xyz$,
$\Omega^1=R^3/R(yz,xz,xy)$ is torsion-free of rank two and is
not free at the triple point.  If $I=(xy,xz,yz)$, there is an exact
sequence
\[
0\to T\to\Omega^2\xrightarrow{\wedge dt} I\to0,
\qquad
T=\frac{R u\oplus R v}{(xu,\ y(u-v),\ zv)},
\quad IT=0,
\]
and $\Omega^3=R/I$.  Finally, for $t=s^m$ in
$R=\C\{s,y_1,y_2\}$,
\[
\begin{aligned}
\Omega^1&=R^2\oplus R/(s^{m-1})\,ds,\\
\Omega^2&=R\,dy_1\wedge dy_2
 \oplus\bigl(R/(s^{m-1})\bigr)^2,\\
\Omega^3&=R/(s^{m-1})\,ds\wedge dy_1\wedge dy_2.
\end{aligned}
\]
\end{claim}
\begin{proof}
The degree-one presentation is the quotient by $dt$; higher degrees
are its exterior powers.  At a double point the transverse quotient
$R^2/R(y,x)$ maps isomorphically to $(x,y)$ by
$(a,b)\mapsto xa-yb$ after an immaterial change of sign of one
basis vector.  More explicitly one may use
$(a,b)\mapsto xa-yb$ with relation $(y,x)$; its kernel is precisely
that relation, since $x,y$ are relatively prime.  In degree two,
with basis $A=dx\wedge dy,B=dx\wedge dz,C=dy\wedge dz$, the
relations are $xA=yA=0$, $yB+xC=0$.  This gives the displayed
torsion summand and the ideal $(x,y)$; wedging once more gives
$R/(x,y)$.

At the triple point the map from ordinary to logarithmic relative
one-forms sends $(a,b,c)$ to $(xa-zc,yb-zc)$.
Its kernel is precisely $R(yz,xz,xy)$: the two equations imply
$c=xyh$, $a=yzh$, $b=xzh$ by relative primality of the coordinates.
Thus it injects into $R^2$, proving torsion-freeness.  Its fibre at
the maximal ideal has dimension three although its rank is two,
so it is not free.  The same minimal-generator argument works at
the double locus.

In degree two, wedging with $dt$ is the row
$(xy,-xz,yz)$ in the basis $(A,B,C)$.
Its syzygies are generated freely by
$u=(z,y,0)$ and $v=(0,y,x)$.
To check generation, an equation $xyA-xzB+yzC=0$ implies
$x\mid C$, say $C=xC_1$; then $yA=z(B-yC_1)$ implies
$A=zA_1$, $B=y(A_1+C_1)$.  Thus the syzygy is $A_1u+C_1v$.
The three relations from $dt\wedge(dx,dy,dz)$ are
$-xu,y(u-v),zv$, proving the stated presentation of the kernel.
The ideal $I$ kills it by those relations.  Conversely the quotient
$I$ is torsion-free, so this kernel is exactly the torsion of
$\Omega^2$.  The presentation immediately gives $\Omega^3=R/I$.
All matrix identities have also been checked over $\Z[x,y,z]$ by
the supplied exact program.  For $t=s^m$ the sole degree-one
relation is $ms^{m-1}ds=0$; $m$ is a unit and taking exterior
powers gives the three formulas.
\end{proof}

\begin{claim}{A4.2}{Refuted}{The claimed degree-one torsion along the cusp}
The assertion in AC26 Remark 9.21, PDF page 81, that
$\Omega_{X/B}^1$ has torsion along $W\cup S_1\cup S_2$ is false.
Its torsion is supported precisely on $S_1\cup S_2$; its degree-one
cusp modules are torsion-free but not locally free on the double
locus.  The plural statement in Remark 1.3, PDF page 5, must be
qualified by degree: $\Omega^2_{X/B}$ does have torsion on the cusp
double locus.  The workbench's degree-one correction is correct;
it must not be extended to degree two.
\end{claim}
\begin{proof}
Both disputed sentences were checked on rendered pages of the
original PDF, so this is not a transcription correction.  The
double and triple charts in A4.1 exhaust the singular strata of
$W$.  They prove absence of degree-one torsion at the cusp and
exhibit nonzero degree-two torsion there.  At a multiple fibre the
finite etale comparison
$q^*\Omega^1_{N_j/B}=\Omega^1_{J'_j/B}$ is an isomorphism,
and its last summand in A4.1 is nonzero for $m=3,4$.
Faithful flatness descends both its existence and its support.
There are no other critical fibres and the ordinary relative
cotangent is locally free on their complement.  Thus the stated
support is exact.  The erroneous sentence is at transcription
line 4330 (the final paragraph of Remark 9.21); the workbench's
correct degree-one replacement occupies lines 980--1006 of
\texttt{hodge\_frolicher.tex}.  No numerical Hodge value in that
workbench is refuted by this calculation.
\end{proof}
